%!TEX root = ../../main.tex
%% ===============================================================
%% 2.5 예측기 비교를 위한 통계 검정
%% 파일: chapters/02_background/05_statistical_tests.tex
%% 상위: chapters/02_background/chapter.tex
%% 정의 라벨: sec:bg-stats
%% 참조 라벨: -
%% 포함 객체: -
%% 인용: delong1988comparing, efron1979bootstrap, holm1979simple, mcnemar1947note
%% 상태: 초안 (docs/OUTLINE.md 참고)
%% ===============================================================

\section{예측기 비교를 위한 통계 검정}
\label{sec:bg-stats}

\begin{definition}[DeLong 검정]
동일한 시험 집합에서 얻은 두 AUC $\hat{A}_1, \hat{A}_2$의 차이를 검정한다. AUC를 $U$ 통계량으로 보고 그 점근 분산·공분산을 구조적 성분(structural component)으로 추정하면, 귀무가설 $A_1 = A_2$ 아래에서
\begin{equation}
  Z = \frac{\hat{A}_1 - \hat{A}_2}{\sqrt{\widehat{\mathrm{Var}}(\hat A_1) + \widehat{\mathrm{Var}}(\hat A_2) - 2\,\widehat{\mathrm{Cov}}(\hat A_1, \hat A_2)}} \;\sim\; \mathcal{N}(0,1)
\end{equation}
이다. DeLong 등 \cite{delong1988comparing}이 제안한 비모수 검정으로, 상관된(같은 표본에 대한) AUC 비교에 적합하다.
\end{definition}

\begin{definition}[McNemar 검정]
두 분류기의 정오표에서 $b$ = A만 맞힌 수, $c$ = B만 맞힌 수라 할 때 연속성 보정된 통계량 $\chi^2 = (|b-c|-1)^2/(b+c)$는 귀무가설 아래 자유도 1의 $\chi^2$ 분포를 따른다 \cite{mcnemar1947note}.
\end{definition}

\begin{definition}[Holm--Bonferroni 보정]
$m$개의 가설에 대한 $p$값을 오름차순 $p_{(1)} \le \dots \le p_{(m)}$으로 정렬하고, $p_{(i)} \le \alpha/(m-i+1)$이 처음으로 실패하는 $i$ 이전까지만 기각한다. 가족 오류율(family-wise error rate)을 $\alpha$로 통제하면서 Bonferroni보다 균일하게 더 검정력이 높다 \cite{holm1979simple}.
\end{definition}

\begin{definition}[부트스트랩 백분위 신뢰구간]
관측 표본으로부터 복원 추출한 $B$개의 재표본 각각에서 통계량을 계산하고, 그 경험 분포의 $2.5$ 및 $97.5$ 백분위를 95\% 신뢰구간으로 취한다. Efron \cite{efron1979bootstrap}이 제안하였다.
\end{definition}
